\section{The integral tensor pair and its four weights}
\label{sec:tensor-integral65}
Fix the source Case III tensor datum $d$. Write $O\subset O_1$ for
the coefficient rings and $K'$ for its auxiliary quadratic field.
The two integral rank-two factors are $\rho_1$ and $\tau$.
For a tower $T$, $Q_T$ has representation
$\delta\,\mathrm{tn}(\rho_1,\tau)|_{G_{T.F}}$ and local matrices
$N_v=\mathrm{flg}(x(P),s(P))$ at $P=T.\mathrm{pl}(v)$.
The inclusion of coefficient rings is understood in products with
$\delta$. Write $\phi=\xi\psi^{\pm1}$ with the sign of the initial
label, and $\Phi=\chi_\ell^{k-1}\phi$. All matrix indices start at zero.

% BEGIN Q65P1
\paragraph{Q65P1. The tensor multiplier and the local change of basis.}
The signed tensor matrix $\mathrm{tn}(G,H)$ represents $G\otimes H$
in the ordered basis
\[
 e_0\otimes f_0,\quad e_0\otimes f_1,\quad
          -e_1\otimes f_1,\quad e_1\otimes f_0 .
\]
Tensor the determinant pairing of the $e$-plane with
$\operatorname{diag}(1,-1)$ on the $f$-plane. Its only nonzero
upper-triangular entries in this basis are the $03$ and $12$ entries,
both one. This is the alternating form defining the source adjoint.
Thus if $H$ preserves the stated symmetric form and
$X=z\,\mathrm{tn}(G,H)$, then $\operatorname{adj}(X)X=z^2\det(G)I_4$. This is the multiplier identity over any commutative
ring: a $2\times2$ matrix pulls back the determinant pairing by its
determinant, without division. In particular it applies integrally
to $G=\rho_1(g)$ and $H=\tau(g)$.

The local matrix, with $x=x(P)$ and $s=s(P)$, is
\[
 N=\begin{pmatrix}
 1&0&0&0\\ s&1&0&0\\-xs&-x&1&0\\ x&0&-s&1
 \end{pmatrix}.
\]
Pairing its columns shows that $N$ preserves the alternating form;
it is lower unitriangular, so its adjoint is its inverse.
Here is the flag calculation used in both conclusions below.
Write $G=\left(\begin{smallmatrix}p&q\\c&r\end{smallmatrix}\right)$,
$H=\left(\begin{smallmatrix}a&b\\b&a\end{smallmatrix}\right)$.
Assume $c+rx=(p+qx)x$ and $s^2=1$, and put
\[
 P_+=p+qx,\qquad P_-=r-xq,\qquad Z_+=a+bs,\qquad Z_-=a-bs .
\]
Multiplying the displayed columns by the signed tensor and expressing
the answers in those same columns gives
\[
 N^{-1}\mathrm{tn}(G,H)N=
 \begin{pmatrix}
 Z_+P_+ & bP_+ &-q(as+b)&aq\\
 0& Z_-P_+ &0&-q(as-b)\\
 0&0& Z_+P_-&-bP_-\\
 0&0&0& Z_-P_-
 \end{pmatrix}.
\]
In particular the first column is an eigenvector. The formula uses
only the two displayed relations, so its triangularity and the
ordered diagonal survive every coefficient map, even in residue
characteristic two.
% END Q65P1

% BEGIN Q65P2
\paragraph{Q65P2. Forbidden inertia and normalization of the integral representation.}
Assume $T_0\le T$, $T.\mathrm{Ok}$, and fix a finite set $W$ of
auxiliary primes. Suppose $\theta$ is trivial on every inertia group
of $G_{T.F}$ away from $\{\ell\}\cup W$, in the sense of
$\mathrm{PairOf4}$. For $h$ in the forbidden set of the frame there
is such a place $P$ with $h$ in its inertia and $\bar r_D(h)=1$.
Take its absolute image $g$. It is an inertia element away from
$\ell$, hence $\chi_\ell(g)=1$. The killing clause for $\delta$
gives $\delta(g)=1$ and $\phi(g)=1$.

We also have $g\in G_{K'}$. In fact $g$ fixes $T_0.F$, since
$T_0\le T$. All four labelled residual diagonal characters are one
at $g$ because $\bar r_D(g)=1$. In particular the first root,
the ratio of the first two such characters, has value one.
If $g$ were outside $G_{K'}$, the off-field clause of
$\mathrm{QuadOn}(T_0,K')$ would make it $-1$. These are distinct
in the source residue field, of odd characteristic. The unramified
clause of $\mathrm{TauDat}$ now gives $\tau(g)=I_2$.
By the hypothesis on $\theta$ and $\phi(g)=1$, the away clause
of $\mathrm{Mod4}$ gives $\rho_1(g)=I_2$. This proves $Q_T.r(h)=I_4$
before any specialization.

By the tower relation the absolute image of $T.g_0$ is an
$\ell^j$-th power of the original marked element for some $j\ge0$.
The original $\rho_1$ there was diagonal, and $\tau$ there was
$\operatorname{diag}(1,-1)$. Powers remain diagonal and their
signed tensor and scalar multiple remain diagonal. This gives the
required normalization, also before specialization.
% END Q65P2

% BEGIN Q65P3
\paragraph{Q65P3. All six clauses of the source pair after arbitrary specialization.}
For $T,W$ and the inertia hypothesis of Q65P2 take \emph{any}
commutative ring $R$ and ring map $f:O_1\to R$. Assume for every
absolute $g$ fixing $T.F$ the source determinant identity
$f(\det\rho_1(g))=f(\Phi(g))$, with $O$ included in $O_1$.
The assertion is $\mathrm{IsPair}$ for $Q_T$ mapped by $f$, with
the frame coefficient map $O\hookrightarrow O_1\xrightarrow{f}R$.

Q65P1 gives multiplier
$f(\delta(g)^2\det\rho_1(g))$. The assumed determinant identity
and $\delta^2\Phi=D.\mathrm{th}$ make this exactly the prescribed
multiplier. Q65P2 gives killing and normalization after applying
$f$. Q65P1 gives lower unitriangular symplectic matrices at every
place. For $g$ in a local decomposition group at $P\mid\ell$, the
splitting clause of $\mathrm{QuadOn}$ puts its absolute image in
$G_{K'}$. Thus $H=\tau(g)$ has the symmetric shape of Q65P1;
$\mathrm{Mod4}$ gives precisely $c+rx=(p+qx)x$. The displayed
conjugation proves the upper-triangular flag condition.

For the final clause let $h\in\mathrm{Mu}_v$. By definition it
lies in the local decomposition group, with an $\ell$-power of
$h$ in the closed commutator subgroup. Its unramified image is zero,
since the local unramified quotient is torsion-free. Its cyclotomic
value has $\ell$-power order in $\mathbb Z_\ell^\times$, hence equals
one: for odd $\ell$ its finite torsion has order dividing $\ell-1$.
Thus $h$ is in inertia with trivial cyclotomic character. The
$\mathrm{ineL}$ clause of $T.\mathrm{Ok}$ makes $\bar r_D(h)=1$,
so $\delta(h)=1$. The ordinary inertia data give
$P_+=\chi_\ell(h)^{k-1}=1$, $P_-=1$ and
$Z_+=\chi_\ell(h)^m=1$. The symmetric form of $H$ gives
$a^2-b^2=1$, hence $Z_+Z_-=a^2-b^2s^2=1$ and $Z_-=1$.
All four diagonal terms of Q65P1 equal one, as do their images
under $f$. This completes the source $\mathrm{PairOf4}$ in its
full determinant and ramification scope.
% END Q65P3

% BEGIN Q65P4
\paragraph{Q65P4. One open subgroup realizes the four prescribed weights.}
For this assertion $T$ can be any tower for $D$. Take an open
subgroup $U\subseteq G_{\mathbb Q}$ and an integer $j$ such that
$\delta^2=\chi_\ell^j$ on $U$. In the source definition of
$T.\mathrm{Wts}$ use the single open subgroup
$(\operatorname{res}_{T.F/\mathbb Q})^{-1}U$.
For every place $P=T.\mathrm{pl}(v)$ and every $h$ there in
inertia, the local splitting of $K'$ still puts its absolute image
in $G_{K'}$. The computation of Q65P1 applies on this tower as
well. Now its inertia values are exactly
\[
 P_+=\chi_\ell^{k-1},\qquad P_-=1,\qquad
 Z_+=\chi_\ell^m,\qquad Z_-=\chi_\ell^{-m}.
\]
The last equality follows again from $Z_+Z_-=1$. After embedding
$O_1$ in $\overline{\mathbb Q}_\ell$, the flag characters are
$\delta$ times $\chi_\ell$ to the respective powers
$k-1+m,k-1-m,m,-m$. On the chosen subgroup their squares are
therefore $\chi_\ell$ to the ordered weights
\[
 j+2(k-1+m),\quad j+2(k-1-m),\quad j+2m,\quad j-2m .
\]
The inequality $2m+3\le k$ makes the source natural-number
subtraction $k-1$ agree with integer subtraction. This proves
$\mathrm{WtsOf4}$, with the specified places, the given integer
$j$, and a single subgroup for all four weights.
% END Q65P4
